\chapter{Introduzione}
\section{Scenario applicativo}
Negli ultimi anni i dispositivi mobili hanno assunto un ruolo centrale nell’ambito dell’informatica personale e professionale.
Tra i sistemi operativi per dispositivi mobili, Android rappresenta una delle piattaforme più diffuse a livello mondiale~\cite{android_marketshare} e 
viene utilizzato in una grande varietà di dispositivi, dagli smartphone ai sistemi embedded. La complessità crescente delle 
applicazioni e del sistema operativo rende sempre più importante disporre di strumenti in grado di analizzare e comprendere il 
comportamento delle applicazioni durante l’esecuzione.

Il monitoraggio delle attività delle applicazioni è fondamentale in diversi ambiti, tra cui il debugging del software, 
l’analisi delle prestazioni e la sicurezza informatica. In particolare, la possibilità di osservare le interazioni tra le 
applicazioni e il sistema operativo consente di individuare anomalie di funzionamento, colli di bottiglia prestazionali e 
potenziali comportamenti malevoli.  

Tradizionalmente il tracing del sistema operativo viene realizzato mediante strumenti specifici che permettono di raccogliere 
informazioni sull’esecuzione dei processi e sulle attività del kernel. Tali strumenti possono tuttavia introdurre un significativo 
overhead computazionale oppure richiedere modifiche al sistema operativo, rendendone difficile l’utilizzo in ambienti reali. 
Inoltre, molti strumenti di analisi disponibili nei sistemi operativi tradizionali, e in particolare in ambiente Linux, non 
sono direttamente utilizzabili su Android senza adattamenti specifici.

Nel 2014 è stato introdotto l’extended Berkeley Packet Filter (eBPF) in Linux~\cite{ebpf_kernel}, una tecnologia che 
permette di eseguire piccoli programmi in modo sicuro ed efficiente all’interno del sistema operativo. Questo consente di osservare 
il comportamento delle applicazioni e del sistema senza la necessità di modifiche al codice sorgente o dello sviluppo di moduli dedicati. 
Per queste ragioni eBPF è diventato uno strumento sempre più utilizzato per attività di tracing e osservabilità nei sistemi Linux~\cite{ebpf_observability}.
\begin{figure}[h]         
    \centering
    \includegraphics[width=0.8\textwidth]{images/architettura-ebpf.png}
    \caption{Architettura di eBPF nel kernel Linux~\cite{ebpf_observability}}
    \label{fig:architettura-ebpf}
\end{figure}

\section{Criticità}
Tra gli strumenti che permettono di utilizzare eBPF in modo pratico, bpftrace costituisce una soluzione particolarmente efficace per 
il tracing delle applicazioni e del sistema operativo.\texttt{ bpftrace}~\cite{bpftrace_docs} mette a disposizione un linguaggio di scripting ad alto livello che 
consente di definire sonde e azioni associate agli eventi osservati, permettendo di raccogliere informazioni durante l’esecuzione senza 
sviluppare programmi eBPF completi. Questo approccio semplifica notevolmente la realizzazione di strumenti di analisi dinamica.

Tuttavia, l’utilizzo di \texttt{ bpftrace} in ambiente Android presenta diverse difficoltà pratiche. In molti dispositivi l’accesso alle 
funzionalità di basso livello del sistema è limitato e spesso non sono disponibili gli strumenti e le librerie necessari per l’esecuzione di programmi basati su eBPF. Inoltre, l’installazione diretta di tali strumenti sul sistema Android può risultare complessa 
o non praticabile, soprattutto al di fuori di ambienti di sviluppo, quando non si ha accesso a certi privilegi che consentono 
di effettuare l'installazione e l'uso di tali strumenti. Ottenere tali 
permessi richiederebbe di effettuare il rooting del dispositivo o di sostituire la release 
Android installata, operazioni spesso non praticabili in contesti reali.

\begin{table}[h]
    \centering
    \begin{tabular}{lcc}
        \toprule
        & \textbf{Emulatore} & \textbf{Dispositivo di Produzione} \\
        \midrule
        \textbf{Accesso root}           & Disponibile     & Richiede rooting   \\
        \textbf{Installazione strumenti} & Possibile       & Richiede rooting  \\
        \textbf{Modifica della ROM}     & Non necessaria  & Spesso necessaria \\
        \textbf{Utilizzo di bpftrace}   & Fattibile       & Non praticabile   \\
        \textbf{Riproducibilità}        & Alta            & Bassa             \\
        \bottomrule
    \end{tabular}
    \caption{Confronto tra ambiente di test e dispositivo di produzione Android.}
    \label{fig:confronto-ambienti}
\end{table}

Queste limitazioni rendono difficile l’impiego delle tecnologie basate su eBPF per l’analisi del comportamento delle applicazioni Android, 
nonostante il loro potenziale per il monitoring del sistema. Risulta quindi necessario definire un ambiente di lavoro che consenta l’utilizzo 
di bpftrace in modo stabile e riproducibile, permettendo allo stesso tempo l’osservazione delle attività delle applicazioni durante l’esecuzione.

Il problema affrontato in questa tesi consiste quindi nel rendere possibile l’utilizzo di bpftrace per il tracing delle applicazioni 
in ambiente Android, individuando una configurazione che consenta di superare le limitazioni tipiche del sistema e di ottenere dati significativi sull’esecuzione delle applicazioni.

\section{Approccio proposto}
Per superare le difficoltà legate all’utilizzo degli strumenti basati su eBPF in ambiente Android, in questo lavoro è stata 
sviluppata una soluzione che permette di eseguire lo strumento bpftrace in un ambiente controllato e riproducibile, mantenendo 
allo stesso tempo la possibilità di osservare il comportamento delle applicazioni Android.

La soluzione proposta si basa sulla creazione di un ambiente di sviluppo virtualizzato in cui sia possibile eseguire Android e 
utilizzare strumenti tipicamente disponibili nei sistemi Linux tradizionali. A questo scopo è stato utilizzato l’emulatore 
Android Cuttlefish~\cite{cuttlefish_docs}, eseguito su sistema operativo Ubuntu, che consente di disporre di un ambiente Android completo e 
configurabile per attività di sviluppo e sperimentazione.

All’interno dell’ambiente Android è stata installata una distribuzione Linux minimale basata su Debian~\cite{debian_distrib}. Questa distribuzione è 
stata utilizzata come ambiente di lavoro per l’esecuzione di bpftrace e degli strumenti necessari al tracing. L’utilizzo di una 
distribuzione Linux separata ha permesso di evitare la compilazione nativa di bpftrace su Android e di semplificare l’
installazione delle dipendenze necessarie.

In questo modo è stato possibile realizzare un ambiente che combina la flessibilità degli strumenti Linux con la possibilità di 
osservare il comportamento di applicazioni eseguite su Android. Su questa base è stato sviluppato un insieme di script e 
strumenti software che permettono di eseguire sessioni di tracing e di raccogliere automaticamente i dati prodotti da bpftrace.
\begin{figure}[h]
      \centering
      \includegraphics[width=0.6\textwidth]{images/architettura-soluzione.png}
      \caption{Architettura della soluzione proposta: Ubuntu, Cuttlefish e ambiente Debian.}
      \label{fig:architettura-soluzione}
\end{figure}

